\chapter{Introduction}

\section{Context and Motivation}
\label{sec:context}

\par Microservices architecture has become the predominant architectural style for building large-scale, distributed applications in both industry and academia. Organizations adopt microservices to achieve independent deployability, fault isolation, technological heterogeneity and team autonomy, qualities that monolithic architectures struggle to provide as systems grow in size and complexity \cite{Newman2015, Richardson2018}. However, the rapid pace of microservices adoption has far outstripped the availability of automated tooling for architectural quality assurance.

\par In practice, teams frequently introduce architectural anti-patterns such as shared databases, cyclic calls between services, or services that are too small to justify their deployment overhead \cite{Taibi2018}. Left undetected, these decisions accumulate as architectural debt and erode the benefits that motivated the move to microservices in the first place, manifesting as cascading failures across service boundaries, deployment friction that forces coordinated releases and maintenance burdens that slow feature delivery. The cost of remediation grows over time, making early, automated detection a pressing need for practitioners who seek feedback on architectural quality without requiring a fully running system.

\section{Problem Statement}
\label{sec:problem-statement}

\par Despite a growing body of research on microservice anti-patterns, the tools available for their detection suffer from several limitations. Existing approaches such as Arcan \cite{Arcelli2017}, MicroART \cite{Granchelli2017} and MSANose \cite{Walker2020} typically operate at a single level of abstraction, either analyzing code-level smells within individual services or examining architectural-level dependency structure across services, but rarely combine both perspectives into a unified analysis. Their detection scope and configuration options vary, and their outputs often consist of abstract warnings without the source code evidence or remediation guidance that developers need to act on the findings. Furthermore, none of the reviewed tools provide a composite metric that quantifies overall architectural quality and tracks it over successive analyses. This dissertation addresses these gaps by designing, implementing and evaluating a tool that bridges intra-service code smell detection with inter-service dependency analysis, covers ten anti-patterns across several architectural dimensions and reports each finding together with the affected services, a source-code snippet as evidence, remediation text and a corresponding deduction from a composite health score.

\section{Objectives}
\label{sec:objectives}

\par The primary objective of this dissertation is to design, implement and evaluate an automated static analysis tool for detecting architectural anti-patterns in Java-based microservice systems. The work pursues five specific objectives. The first is to implement detectors for ten microservice anti-patterns, spanning four architectural dimensions: service design, communication, data management, and deployment and coupling. The second is to implement a multi-level analysis pipeline that combines intra-service code smell detection, performed by integrating DesigniteJava \cite{Sharma2024}, with inter-service dependency analysis based on Spoon's Abstract Syntax Tree traversal \cite{Pawlak2016} and graph algorithms such as Tarjan's strongly connected components algorithm \cite{Tarjan1972}. The third objective is to develop a composite health score that quantifies overall architectural quality on a 0-100 scale with letter grading, and decomposes it into four interpretable categories so that teams can identify which quality dimension is degraded. Fourth, the tool is to be delivered as a web application with a RESTful API, supporting both interactive use through a browser-based dashboard and programmatic integration into CI/CD pipelines. Fifth and finally, the tool is to be evaluated on a corpus of open-source microservice projects to characterize its detection behaviour across anti-pattern types and project scales, with its capabilities positioned against those of related work.

\section{Contributions}
\label{sec:contributions}

\par This dissertation makes five main contributions to the field of automated architectural quality assurance for microservice systems. The first contribution is a multi-level static analysis approach that bridges code-level analysis and architectural-level anti-pattern detection, using intra-service structural metrics such as class cohesion as signals for architectural-level problems such as the God Service anti-pattern. The second contribution is detection support for ten anti-patterns across four architectural dimensions, with configurable thresholds and evidence-based reporting that attaches source code snippets to each finding, enabling developers to understand and remediate issues without additional manual inspection. Third, the work introduces an automated microservice boundary detection mechanism that identifies services from build files (\texttt{pom.xml}, \texttt{build.gradle}) and filters out non-service modules, without the running system or the manual curation of service boundaries that existing tools tend to rely on. The fourth contribution is a composite health score with category-based decomposition and letter grading that enables teams to track architectural quality over time and measure the concrete impact of refactoring efforts through a built-in analysis diff feature. Finally, the fifth contribution is the MSA Detector itself, an open-source web application that implements all of the above, supporting both ZIP archive upload and Git repository cloning as project ingestion modes.

\section{Thesis Structure}
\label{sec:thesis-structure}

\par The remainder of this dissertation is organized as follows. Chapter~\ref{chap:ch2} reviews the theoretical background on microservices architecture and anti-patterns, surveys related detection tools and approaches, and positions this work within the existing landscape. Chapter~\ref{chap:ch3} presents the detection methodology, describing the analysis pipeline, the dependency graph construction and the detailed detection algorithms for each anti-pattern. Chapter~\ref{chap:ch4} describes the system architecture, data model, backend and frontend implementation, and deployment configuration. Chapter~\ref{chap:ch5} presents the evaluation datasets, detection results, manual validation of the findings, and discusses limitations and threats to validity. Finally, Chapter~\ref{conclusions} summarizes the contributions, reflects on the implications of the results and outlines directions for future research.

\section{Use of Generative AI Tools}
\label{sec:genai}

\par In the interest of transparency, and in accordance with the academic integrity guidelines governing the use of generative artificial intelligence in student work, this section discloses where and how such tools were used in the preparation of this dissertation. Generative AI assistants, namely Anthropic's Claude (accessed through the Claude Code command-line interface), were employed in a limited and supervised capacity. The author retains full responsibility for the entire content of this work, for all design and implementation decisions, and for the results reported herein. Every output produced with the assistance of these tools was reviewed critically, verified, and, where necessary, corrected or rewritten before being incorporated, and the tools were used throughout with discernment and academic integrity.

\par In the writing of the dissertation, generative AI tools were used to improve the readability of the text, to suggest alternative phrasings and reformulations of passages first drafted by the author, and to flag wording that read as artificial or repetitive so that it could be revised. The intellectual content, argumentation, structure, and conclusions of the dissertation are the author's own, and no portion of the text was adopted from AI-generated output without revision and verification.

\par In the development of the accompanying software, generative AI tools were used to assist with routine and repetitive coding tasks, such as generating boilerplate and scaffolding, and to provide review feedback on correctness, simplification, and potential edge cases. All source code was read, understood, tested, and integrated by the author, who remains responsible for its correctness and behaviour.
